\usepackage{fontspec}
\usepackage{xeCJK}
\setmainfont{Latin Modern Roman}
\setCJKmainfont[ItalicFont={Noto Serif CJK JP},BoldItalicFont={Noto Serif CJK JP Bold}]{Noto Serif CJK JP}
\setCJKsansfont{Noto Sans CJK JP}
\setCJKmonofont{Noto Sans Mono CJK JP}
\newfontfamily\proseunicodefont{DejaVu Serif}
\usepackage[margin=1in]{geometry}
\usepackage{microtype}
\usepackage{amsmath,amssymb,amsthm,mathtools,bm}
\usepackage{graphicx}
\usepackage{xcolor}
\usepackage{booktabs}
\usepackage{array}
\usepackage{longtable}
\usepackage{enumitem}
\usepackage{float}
\usepackage{caption}
\usepackage{hyperref}
\usepackage{fancyhdr}
\usepackage{tikz}
\usetikzlibrary{arrows.meta,positioning,calc,decorations.pathmorphing}

\hypersetup{
  colorlinks=true,
  linkcolor=blue!50!black,
  urlcolor=blue!50!black,
  citecolor=blue!50!black
}

\setlength{\parskip}{0.5em}
\setlength{\parindent}{0pt}
\setlist[itemize]{topsep=0.25em,itemsep=0.2em}
\setlist[enumerate]{topsep=0.25em,itemsep=0.2em}
\setlength{\headheight}{40pt}
\numberwithin{equation}{chapter}

\newtheorem{theorem}{定理}[chapter]
\newtheorem{proposition}[theorem]{命題}
\newtheorem{lemma}[theorem]{補題}
\newtheorem{corollary}[theorem]{系}
\theoremstyle{definition}
\newtheorem{definition}[theorem]{定義}
\newtheorem{example}[theorem]{例}
\theoremstyle{remark}
\newtheorem{remark}[theorem]{注意}

\providecommand{\ket}[1]{\left\lvert#1\right\rangle}
\providecommand{\bra}[1]{\left\langle#1\right\rvert}
\providecommand{\braket}[2]{\left\langle#1\,\middle|\,#2\right\rangle}
\providecommand{\expect}[1]{\left\langle#1\right\rangle}
\newcommand{\lectureoverview}[1]{\paragraph{概観。} #1}
\newenvironment{classroomqa}{%
  \par\medskip
  \begin{minipage}{0.96\textwidth}
}{%
  \end{minipage}
  \par\medskip
}
\newcommand{\audiencequestion}[1]{\noindent\textbf{受講者からの質問。} #1\par\smallskip}
\newcommand{\lecturerresponse}[1]{\noindent\textbf{回答。} #1\par}
\newcommand{\lecturetimestamp}[1]{\footnote{講義時刻： #1.}}
\newcommand{\lectureframe}[1]{\par\smallskip{\centering\footnotesize\itshape 講義映像： #1.\par}}
\newcommand{\editorialnote}[1]{\footnote{編集上の補足： #1}}

\makeatletter
\newcommand{\course@headerfont}{\normalfont\small\itshape}
\newcommand{\course@pagefont}{\normalfont\small}
\newcommand{\course@headerraise}{14pt}
\newlength{\course@headerboxheight}
\setlength{\course@headerboxheight}{\dimexpr\headheight-\course@headerraise-9pt\relax}
\newcommand{\course@leftheadbox}[1]{%
  \raisebox{\course@headerraise}{%
    \parbox[t][\course@headerboxheight][t]{0.82\textwidth}{%
      \raggedright\course@headerfont\strut\MakeUppercase{#1}\strut%
    }%
  }%
}
\newcommand{\course@rightheadbox}[1]{%
  \raisebox{\course@headerraise}{%
    \parbox[t][\course@headerboxheight][t]{0.82\textwidth}{%
      \raggedleft\course@headerfont\strut\MakeUppercase{#1}\strut%
    }%
  }%
}
\renewcommand{\chaptermark}[1]{\markboth{\MakeUppercase{\chaptername\ \thechapter.\ #1}}{}}
\renewcommand{\sectionmark}[1]{\markright{\MakeUppercase{\thesection\hspace{0.4em}#1}}}
\fancyhf{}
\fancyhead[L]{\course@leftheadbox{\leftmark}}
\fancyhead[R]{\course@pagefont\thepage}
\renewcommand{\headrulewidth}{0.5pt}
\pagestyle{fancy}
\fancypagestyle{plain}{%
  \fancyhf{}
  \fancyhead[L]{\course@leftheadbox{\leftmark}}
  \fancyhead[R]{\course@pagefont\thepage}
  \renewcommand{\headrulewidth}{0.5pt}
}
\makeatother

\renewcommand{\contentsname}{目次}
\renewcommand{\chaptername}{章}
\renewcommand{\figurename}{図}
\renewcommand{\tablename}{表}
\renewcommand{\proofname}{証明}
\renewcommand{\bibname}{参考文献}
\emergencystretch=2em
\pdfstringdefDisableCommands{\def\proseunicodefont{}}

\geometry{paperwidth=177.6mm,paperheight=252mm,left=16mm,right=16mm,
 top=18mm,bottom=18mm,includeheadfoot,headheight=28pt,headsep=7mm,footskip=9mm}
\usepackage{scrextend}
\changefontsizes[17pt]{12pt}
\raggedbottom
\renewenvironment{classroomqa}{\par\medskip\begingroup}{\par\endgroup\medskip}
\fancyhf{}
\setlength{\headwidth}{\textwidth}
\fancyhead[L]{\parbox{\textwidth}{\raggedright\fontsize{9}{11}\selectfont\nouppercase{\leftmark}}}
\fancyfoot[C]{\small\thepage}
\renewcommand{\headrulewidth}{0.3pt}
\fancypagestyle{plain}{\fancyhf{}\fancyfoot[C]{\small\thepage}\renewcommand{\headrulewidth}{0pt}}
